\section{PRESERVATION OF ARGUMENT: SATISFACTION OF REGULATORY CRITERIA MANDATES APPROVAL; THE ``FINAL MERITS DETERMINATION'' IS UNLAWFUL}
\label{sec:mukherji}

The beneficiary respectfully preserves the following legal argument for the record, citing recent federal court precedent that directly addresses the lawfulness of USCIS's current adjudication framework.

\subsection{The Regulatory Standard and Statutory Command}

The Immigration and Nationality Act establishes that an alien ``who seeks to enter the United States to continue work in the area of extraordinary ability'' shall be granted classification if the alien ``has extraordinary ability in the sciences, arts, education, business, or athletics which has been demonstrated by sustained national or international acclaim.'' 8 U.S.C. § 1153(b)(1)(A)(i).

The implementing regulation at 8 C.F.R. § 204.5(h)(3) provides that evidence of extraordinary ability shall consist of either (1) a one-time achievement of a major, internationally recognized award, or (2) ``at least three of the following'' ten enumerated criteria. The regulation does not state that meeting three criteria is merely a threshold for further discretionary review. The plain text establishes that satisfaction of three criteria constitutes sufficient evidence to demonstrate extraordinary ability.

\subsection{The Unlawful``Final Merits Determination'' Framework}

In \textit{Kazarian v. USCIS}, 596 F.3d 1115 (9th Cir. 2010), the Ninth Circuit outlined a two-step process: first, determine whether the beneficiary has provided qualifying evidence under at least three criteria; second, evaluate the totality of the evidence to determine whether the beneficiary has demonstrated sustained national or international acclaim. This framework was adopted by USCIS through the Administrative Appeals Office and incorporated into policy through internal memoranda—most notably the December 22, 2010 Policy Memorandum PM-602-0005.1.

Critically, this ``final merits determination'' was \textbf{never promulgated through notice-and-comment rulemaking} as required by the Administrative Procedure Act, 5 U.S.C. § 553. The APA mandates that substantive rules affecting individual rights must be adopted through public notice and opportunity for comment. USCIS implemented the final merits framework through policy guidance, bypassing this procedural requirement entirely.

\subsection{The \textit{Mukherji v. Miller} Decision}

On January 28, 2026, the United States District Court for the District of Nebraska issued a landmark decision in \textit{Mukherji v. Miller}, Case No. 4:24-cv-03170 (D. Neb. Jan. 28, 2026). In that case, USCIS conceded that the petitioner, Anahita Mukherji, an Indian journalist, met at least five of the ten regulatory criteria—exceeding the statutory minimum of three. Despite this concession, USCIS denied the petition at the ``final merits determination'' stage, concluding she had not demonstrated sustained national or international acclaim.

Senior Judge Joseph F. Bataillon ruled that USCIS's denial was \textbf{arbitrary, capricious, and contrary to law}. The court found that the ``final merits determination'' framework was \textbf{unlawfully adopted} without compliance with APA rulemaking procedures. The court vacated the denial and \textbf{remanded with instructions to approve the petition}, holding that once the regulatory criteria are satisfied, denial based on an extra-regulatory, subjectively-applied standard violates the APA.

The court's reasoning is dispositive: \textbf{USCIS may not impose requirements beyond those codified in the regulations through properly-adopted rulemaking procedures}. The final merits determination, as currently applied, functions as an unlawful substantive rule that was never subjected to the notice-and-comment process mandated by 5 U.S.C. § 553. When an agency bypasses APA rulemaking requirements, the resulting policy is void and may not be applied to deny benefits. \textit{See Perez v. Mortgage Bankers Ass'n}, 575 U.S. 92, 101 (2015) (agencies must use notice-and-comment procedures when adopting substantive rules); \textit{Chrysler Corp. v. Brown}, 441 U.S. 281, 313 (1979) (agency action that imposes new obligations without APA compliance is invalid).

Crucially, the \textit{Mukherji} court did not merely find the denial improper—it \textbf{ordered USCIS to approve the petition}. This remedy confirms that once the regulatory criteria are satisfied, approval is \textbf{mandatory}, not discretionary, absent a validly-adopted rule imposing additional requirements.

\subsection{Application to This Petition}

The beneficiary in this case has demonstrably met \textbf{at least six of the ten regulatory criteria}:

\begin{enumerate}
    \item \textbf{Original Contributions of Major Significance} (8 C.F.R. § 204.5(h)(3)(v)): Machine learning methodologies permanently archived in the U.S. National Transportation Library, adopted by the Florida Department of Transportation, and deployed in production systems serving 80+ million users at JPMorgan Chase.
    
    \item \textbf{Scholarly Articles} (8 C.F.R. § 204.5(h)(3)(vi)): Eight peer-reviewed publications in professional journals, with federal government archival and cross-state adoption.
    
    \item \textbf{Judging the Work of Others} (8 C.F.R. § 204.5(h)(3)(iv)): Reviewer for the NeurIPS 2025 Position Paper Track, the ICLR 2025 FM-Wild Workshop, and the ICML 2026 AIWILD Workshop---specialized tracks and workshops within the premier machine learning conferences globally---and service on the IEEE Senior Member Application Review Panel.
    
    \item \textbf{Membership in Associations Requiring Outstanding Achievements} (8 C.F.R. § 204.5(h)(3)(ii)): IEEE Senior Member, a grade requiring ten years of professional experience and five years of significant performance as judged by existing Senior Members or Fellows.
    
    \item \textbf{Published Material About the Alien} (8 C.F.R. § 204.5(h)(3)(iii)): Coverage in IEEE official channels (IEEE Transmitter, IEEE.tv, IEEE Twitter with 342,900 followers), India's largest newspapers (Dainik Bhaskar with 63 million readers, Rajasthan Patrika with 25.9 million readers), and major technical trade publications (TechTarget with 1.4M+ daily readers).
    
    \item \textbf{Leading or Critical Role} (8 C.F.R. § 204.5(h)(3)(viii)): Founding engineer on the Chase Shopping Hub at JPMorgan Chase \& Co., designing backend architecture for a mission-critical platform processing 5+ million monthly requests for 80+ million cardmembers.
\end{enumerate}

Under the plain text of 8 C.F.R. § 204.5(h)(3), satisfaction of three criteria constitutes sufficient evidence. The beneficiary has met six—double the regulatory requirement.

\subsection{The Regulatory Text Is Unambiguous and Self-Executing}

The regulation at 8 C.F.R. § 204.5(h)(3) states that ``[a] petition for an alien of extraordinary ability must be accompanied by evidence that the alien has sustained national or international acclaim and that the alien's achievements have been recognized in the field of expertise. Such evidence \textbf{shall} include evidence of a one-time achievement [or] documentation of at least three of the following''—followed by ten enumerated criteria. (emphasis added).

The word ``shall'' is mandatory, not permissive. \textit{See Kingdomware Techs., Inc. v. United States}, 579 U.S. 162, 172 (2016) (``shall'' imposes nondiscretionary duty). The regulation does not state that meeting three criteria is a threshold for further discretionary analysis. Rather, evidence of three qualifying criteria \textbf{shall constitute} the required showing.

USCIS's practice of denying petitions after conceding that three criteria are met—by applying unpromulgated standards such as citation thresholds, subjective assessments of ``international'' recognition, or comparison to hypothetical ``top researchers'' in unrelated subfields—directly contradicts the regulatory text and constitutes ultra vires agency action.

\subsection{Preservation of Argument Regarding Citation Metrics and Extra-Regulatory Standards}

The beneficiary anticipates that USCIS may attempt to deny this petition at the ``final merits determination'' stage by asserting that the beneficiary's 37 citations are ``insufficient'' compared to academic researchers in theoretical machine learning. Such a denial would be \textbf{ultra vires} for the following reasons:

\textbf{First}, the regulation does not impose a citation threshold. The scholarly articles criterion at 8 C.F.R. § 204.5(h)(3)(vi) requires only ``evidence of authorship of scholarly articles in professional or major trade publications.'' It does not specify a minimum citation count. Imposing such a requirement through the final merits determination constitutes impermissible substantive rulemaking without APA compliance.

\textbf{Second}, the beneficiary's field is \textbf{applied AI/ML systems development for critical infrastructure}, not academic theoretical research. In this field, impact is measured by \textbf{institutional reliance and production deployment}, not citation counts. The beneficiary's work has been:
\begin{itemize}
    \item Permanently archived by the U.S. Department of Transportation in the National Transportation Library
    \item Adopted by the Florida Department of Transportation in federally-funded research
    \item Deployed in production at JPMorgan Chase, a systemically important financial institution, serving 80+ million users
    \item Formally integrated into academic curricula at government-regulated institutions
\end{itemize}

These forms of validation demonstrate extraordinary ability in the beneficiary's actual field—not a hypothetical field of theoretical ML research where citation-counting may be more appropriate.

\textbf{Third}, as \textit{Mukherji} makes clear, USCIS may not deny a petition based on subjective assessments that contradict the regulatory criteria already satisfied. If USCIS concedes that the beneficiary meets the scholarly articles criterion (through publication in peer-reviewed journals), it cannot then deny the petition by imposing an extra-regulatory citation standard at Step 2.

\subsection{Jurisdictional Note and Persuasive Authority}

The beneficiary acknowledges that \textit{Mukherji v. Miller} was decided by the U.S. District Court for the District of Nebraska, which sits within the Eighth Circuit, and that this petition is being filed from Delaware, which sits within the Third Circuit. As such, \textit{Mukherji} is not binding precedent on USCIS adjudicators in the Third Circuit.

However, the legal issue presented by \textit{Mukherji}—whether USCIS violated the APA by implementing substantive requirements without notice-and-comment rulemaking—is a \textbf{question of federal administrative law, not circuit-specific precedent}. The APA is a statute of nationwide application. The holding that agencies may not bypass 5 U.S.C. § 553 to impose new obligations applies uniformly across all circuits. \textit{See Motor Vehicle Mfrs. Ass'n v. State Farm Mut. Auto. Ins. Co.}, 463 U.S. 29, 42 (1983) (agency action that fails to comply with APA procedures is arbitrary and capricious).

Moreover, USCIS is a \textbf{single federal agency} bound by the APA nationwide. If the final merits determination violates the APA in the Eighth Circuit, it violates the APA everywhere. Allowing USCIS to continue applying an unlawful framework in circuits where it has not yet been challenged would create the perverse incentive for agencies to ignore adverse court decisions until compelled to comply circuit-by-circuit—a result inconsistent with the uniform application of federal administrative law.

The beneficiary respectfully urges USCIS to apply \textit{Mukherji}'s reasoning as persuasive authority grounded in the plain requirements of the APA and to approve this petition based on satisfaction of the regulatory criteria, thereby avoiding the need for judicial intervention.

\subsection{Conclusion}

The beneficiary has met \textbf{six of the ten regulatory criteria}—twice the statutory minimum of three. Under the plain text of 8 C.F.R. § 204.5(h)(3), which uses the mandatory term ``shall,'' this constitutes sufficient evidence of extraordinary ability. The regulation does not authorize USCIS to deny petitions based on unpromulgated, subjective assessments of ``international acclaim,'' citation thresholds, or comparisons to researchers in unrelated fields. Such denials violate the APA as established in \textit{Mukherji v. Miller}.

The beneficiary respectfully requests that USCIS approve this petition based on satisfaction of the regulatory criteria, consistent with the statute (8 U.S.C. § 1153(b)(1)(A)), the properly-adopted regulations (8 C.F.R. § 204.5(h)(3)), the Administrative Procedure Act (5 U.S.C. §§ 553, 706), and recent federal court precedent.

\textbf{Any denial based on extra-regulatory standards not subjected to notice-and-comment rulemaking would be arbitrary, capricious, an abuse of discretion, and contrary to law.}

This argument is preserved for the record and for any subsequent administrative or judicial review.
